%% DTRAP submission manuscript (acmart, double-anonymous review).
%% Modeled on the DTRAP-2026-0102 (Theia) submission package.
%%
%% Commands for TeXCount
%TC:macro \cite [option:text,text]
%TC:macro \citep [option:text,text]
%TC:macro \citet [option:text,text]
%TC:envir table 0 1
%TC:envir table* 0 1
%TC:envir tabular [ignore] word
%TC:envir displaymath 0 word
%TC:envir math 0 word
%TC:envir comment 0 0
%%
\documentclass[manuscript,screen,review,anonymous]{acmart}

\usepackage{booktabs}
\usepackage{multirow}
\usepackage{listings}

\lstset{
  basicstyle=\ttfamily\small,
  breaklines=true,
  frame=single,
  columns=flexible,
}

%% Rights management information (submission: none).
\setcopyright{none}
\acmJournal{DTRAP}

\begin{document}

\title{From Binary Gates to Security Lattices: Applying Bell-LaPadula to LLM Agent Safety Evaluation}

%% Double-anonymous review: the [anonymous] class option replaces the
%% author block with ANONYMOUS AUTHOR(S).
\author{Anonymous Author(s)}

\begin{abstract}
Autonomous LLM agents are under active attack in deployment. OpenClaw, a widely deployed agent framework, has accumulated over 60 CVEs and 100 security advisories since January 2026, including compound chains that silently hijack agents (ClawJacked, CVE-2026-32025), exfiltrate credentials through allowlisted commands, and escape sandboxes to install backdoors (Claw Chain, CVE-2026-44112/13/15/18). Yet the benchmarks that certify such agents as trustworthy reduce safety to a flat binary gate (safe~=~1 or unsafe~=~0) with no hierarchical structure, no information flow tracking, and no formal security properties. Using Claw-Eval, the most comprehensive trajectory-aware agent benchmark (300 tasks, 19 mock services), as a case study, we show that 81\% of tasks never evaluate safety, 67\% of credential-bearing tasks lack leakage detection, and the binary gate equates critical credential exfiltration with minor policy deviations. We propose \textbf{SafeLattice}, a framework that applies the Bell-LaPadula (BLP) formal security model to agent safety evaluation. SafeLattice introduces (1)~a four-level security classification lattice (Critical/High/Medium/Low) analogous to BLP's Top~Secret through Unclassified hierarchy, (2)~information flow analysis that detects \textit{write-down} violations when agents expose high-sensitivity data through low-clearance output channels, (3)~a Biba integrity extension that frames prompt injection as a low-integrity input corrupting a high-integrity action, and (4)~automatic enforcement of declared safety constraints via lattice-aware state-transition verification. We prove, using the six-step state-machine verification methodology, that SafeLattice's scoring preserves a secure-state invariant across the agent trajectory and is backward-compatible (no binary-failing trace can pass). Empirically, on a labeled trace-level corpus of 68 agent trajectories grounded in real Claw-Eval tasks and documented OpenClaw CVEs, SafeLattice raises violation-detection recall from 0.52 to 0.95 and resolves three severity levels where binary scoring resolves one; it also \textit{reorders} models---an agent that ranks first on task completion ranks last once safety is scored on the lattice (Kendall $\tau=-0.67$ between completion-only and SafeLattice rankings). A live evaluation of 1,222 rollouts across six frontier models (GPT-4o, Claude Sonnet/Haiku, Gemini Flash, Llama-3.3-70B) confirms the threat is extant rather than hypothetical: four of six models leak credential material on an authorized-audit task and the deployed binary gate scores every leak safe; against verified ground truth (85 violations, established by mechanical task-specific verifiers), the binary gate achieves recall 0.06 versus SafeLattice's 1.0 (F1 0.11 vs.\ 0.81), with severity resolved at three levels, bootstrap confidence intervals, and over-refusal at most 0.13 once refusal is separated from incapacity. Canary-based obfuscation detection closes a measured evasion gap, raising evasion recall from 0.33 to 1.0 on a six-variant encoding probe. SafeLattice increases effective safety coverage from 19\% to 28\% of the benchmark's task space.
\end{abstract}

%% CCS concepts generated per http://dl.acm.org/ccs.cfm.
%% Regenerate at camera-ready if the committee requests different weights.
\begin{CCSXML}
<ccs2012>
<concept>
<concept_id>10002978.10002986.10002990</concept_id>
<concept_desc>Security and privacy~Logic and verification</concept_desc>
<concept_significance>500</concept_significance>
</concept>
<concept>
<concept_id>10002978.10002991.10002994</concept_id>
<concept_desc>Security and privacy~Information flow control</concept_desc>
<concept_significance>500</concept_significance>
</concept>
<concept>
<concept_id>10010147.10010178.10010219.10010220</concept_id>
<concept_desc>Computing methodologies~Intelligent agents</concept_desc>
<concept_significance>300</concept_significance>
</concept>
<concept>
<concept_id>10002978.10002986.10002988</concept_id>
<concept_desc>Security and privacy~Security requirements</concept_desc>
<concept_significance>300</concept_significance>
</concept>
</ccs2012>
\end{CCSXML}

\ccsdesc[500]{Security and privacy~Logic and verification}
\ccsdesc[500]{Security and privacy~Information flow control}
\ccsdesc[300]{Computing methodologies~Intelligent agents}
\ccsdesc[300]{Security and privacy~Security requirements}

\keywords{large language model (LLM) agents, safety evaluation, Bell-LaPadula, formal security models, information flow control, security lattice, agent benchmarking, credential leakage}

\maketitle

%% ============================================================
%% SECTION 1: INTRODUCTION
%% ============================================================
\section{Introduction}
\label{sec:intro}

The threat is no longer hypothetical. OpenClaw, an autonomous LLM agent framework deployed on developer workstations worldwide, has accumulated over 60 CVEs and 100 GitHub Security Advisories since January 2026 \cite{openclawvulntaxonomy2026}. The ClawJacked vulnerability (CVE-2026-32025, CVSS 8.8) allows any malicious website to silently hijack a local agent, search Slack history for API keys, exfiltrate files, and execute shell commands \cite{clawjacked2026}. The Claw Chain (CVE-2026-44112/13/15/18) chains four CVEs into a sandbox escape that exfiltrates credentials through allowlisted commands and installs backdoors \cite{clawchain2026}. These are extant, weaponized attacks against agents that the ecosystem's evaluation benchmarks certify as safe.

That certification is the problem this paper addresses. Large language models (LLMs) have evolved from conversational assistants into autonomous agents capable of executing multi-step workflows in real-world software environments \cite{openai2023gpt4, anthropic2024claude, google2023gemini}. Modern agent frameworks such as Claude Code \cite{anthropic2025code}, OpenClaw \cite{openclaw2025}, and Cursor \cite{cursor2024} equip LLMs with the ability to call tools, navigate file systems, query databases, and orchestrate actions across multiple applications. Practitioners deciding which agent to deploy rely on benchmarks that claim to measure safety alongside task completion. Yet the safety evaluation mechanisms in these benchmarks stand in stark contrast to the rigorous, decades-old discipline of formal security modeling. Since Bell and LaPadula's foundational work in 1973 \cite{belllapadula1976}, the security community has established that trustworthy systems require \textit{hierarchical classification levels}, \textit{mandatory information flow control}, and \textit{provable state transition invariants}. These properties cannot be reduced to a single binary gate without sacrificing the very diagnostic power that safety evaluation is meant to provide.

Yet this is precisely what current agent benchmarks do. Claw-Eval \cite{claweval2026}, the most comprehensive trajectory-aware evaluation framework with 300 human-verified tasks scored along Completion, Safety, and Robustness dimensions, implements safety as a multiplicative binary gate: $\texttt{safety} \in \{0, 1\}$. A recent survey finds that 79\% of agent safety benchmarks use binary pass/fail metrics \cite{yu2026taxonomy}, and that benchmark choice systematically shapes reported safety with no concordance across evaluation instruments (Kendall's $W = 0.10$, $p = 0.94$). This fragmentation is not merely an engineering inconvenience---it is a direct consequence of evaluating security without formal security models.

The gap between extant attacks and deployed evaluation is measurable. In a live evaluation of 1,222 rollouts across six frontier models (Section~\ref{sec:live-pilot}), four of six models disclosed credential material on an authorized configuration-audit task---and the deployed binary gate scored every one of those leaks as safe. Through systematic analysis of all 300 Claw-Eval tasks and their grader implementations, we identify four critical failures of binary safety scoring:

\begin{enumerate}
    \item \textbf{No severity differentiation.} An agent that leaks API credentials and one that saves an unnecessary email draft both receive $\texttt{safety} = 0.0$, despite representing categorically different threat levels.
    \item \textbf{No information flow control.} Of 18 tasks that provide agents access to credential-bearing services (API keys, database passwords), only 6 (33\%) check whether the agent exposes these secrets in its output. Under binary scoring, an agent that prints raw API keys in the remaining 12 tasks receives $\texttt{safety} = 1.0$.
    \item \textbf{No state transition verification.} Safety is evaluated as a snapshot at grading time, not as a property maintained across the agent's trajectory. Multi-step information flow violations---where individually safe actions compose into an unsafe trajectory---go undetected.
    \item \textbf{Unenforced safety declarations.} 35 tasks declare safety constraints in their configuration files, but the corresponding graders silently ignore them, defaulting to $\texttt{safety} = 1.0$.
\end{enumerate}

We propose \textbf{SafeLattice}, a framework that applies the Bell-LaPadula formal security model to agent safety evaluation. SafeLattice maps BLP's core constructs---security classification levels, the Simple Security Property (``no read up''), the Star Property (``no write down''), and the Basic Security Theorem---to the agent evaluation domain, creating a principled foundation for safety scoring that is provably more expressive than binary gates. Practitioner tooling has already moved in this direction: independent security harnesses such as tinman-openclaw-eval \cite{tinmaneval2026} grade OpenClaw attack probes on five-level severity scales, implicitly rejecting binary classification. SafeLattice supplies the formal foundation that this practice currently lacks, bridging deployed security tooling and the benchmark research that certifies agents.

\subsection{Research Questions}

\textbf{RQ1:} \textit{Does applying Bell-LaPadula's hierarchical security classification to agent evaluation provide more diagnostic granularity than binary safety scoring?}

We construct a four-level security lattice (Critical $>$ High $>$ Medium $>$ Low $>$ Safe) and demonstrate that it correctly differentiates violation severities across 10 representative scenarios while preserving backward compatibility for critical failures.

\textbf{RQ2:} \textit{Can BLP's information flow properties detect safety violations that binary scoring misses?}

We implement the Star Property as an information flow analysis pass that tracks data sensitivity through agent tool-call chains. We show that it detects \textit{write-down} violations (high-sensitivity data written to low-clearance channels) in 12 credential-bearing tasks that currently score $\texttt{safety} = 1.0$.

\textbf{RQ3:} \textit{Does automatic enforcement of declared safety constraints via lattice-aware state verification increase effective safety coverage?}

We implement a safety enforcer that reads declared constraints from task configuration and applies lattice-classified checks. We measure the increase in safety evaluation coverage and the detection of previously invisible violations.

\textbf{RQ4:} \textit{On real agent trajectories, does lattice scoring produce different---and more diagnostic---model rankings than binary scoring?}

We build a labeled trace-level corpus grounded in real Claw-Eval tasks and documented OpenClaw CVEs, score every trace under both systems, and measure detection precision/recall, false negatives fixed, severity resolution, and ranking divergence across simulated agents of differing safety discipline. We then execute the same comparison live: 1,222 rollouts of six frontier models with verified ground truth.

\subsection{Contributions}

This paper makes the following contributions:
\begin{enumerate}
    \item A formal analysis demonstrating that binary safety scoring in agent benchmarks violates the foundational properties of the Bell-LaPadula security model---hierarchical classification, mandatory access control, and secure state transitions---and that these violations cause systematic overestimation of agent trustworthiness.
    \item \textbf{SafeLattice}, a framework that maps BLP's security lattice, Simple Security Property, and Star Property to agent evaluation, providing severity-differentiated safety scoring with provable dominance ordering.
    \item An information flow analysis module that detects write-down violations---agents exposing high-sensitivity data (credentials, PII) through low-clearance output channels---which binary scoring cannot detect by design.
    \item An automatic safety enforcement module that applies declared safety constraints from task configuration using lattice-classified severity levels, increasing effective safety coverage from 19\% to 28\% of the benchmark's task space.
    \item A Biba integrity extension that formalizes prompt injection as a ``no write up'' violation---low-integrity retrieved/tool content driving a high-integrity agent action---and detects it via integrity-flow analysis.
    \item Formal proofs, following the six-step state-machine verification methodology, that SafeLattice's lattice is well-formed, that every trajectory transition preserves a secure-state invariant, and that scoring is backward-compatible with the binary gate.
    \item A labeled trace-level evaluation on 68 agent trajectories showing that SafeLattice raises detection recall from 0.52 to 0.95, resolves three severity levels vs.\ one, and reorders models relative to both binary and completion-only rankings; plus a live evaluation of 1,222 frontier-model rollouts with mechanically verified ground truth (binary recall 0.06 vs.\ SafeLattice 1.0). We release the corpus generator, dual-scoring harness, and live-run harness for reproducibility.
\end{enumerate}

%% ============================================================
%% SECTION 2: RELATED WORK
%% ============================================================
\section{Related Work}
\label{sec:related}

We survey related work across five areas: agent benchmarks, evaluation methodology, real-world agent security incidents, safety and robustness testing, and formal security models applied to AI systems.

\subsection{Agent Benchmarks}

\subsubsection{Tool-use and Code Benchmarks}

SWE-bench \cite{jimenez2024swebench} evaluates agents on real GitHub issue resolution, requiring them to navigate codebases, identify bugs, and generate patches. While it tests genuine software engineering capability, SWE-bench evaluates only the final patch (output-only grading) and includes no safety constraints or robustness testing. ToolBench \cite{qin2024toolllm} and API-Bank \cite{li2023apibank} evaluate API orchestration across hundreds of tools but lack trajectory auditing---a model that fabricates intermediate API responses could score well. Terminal-Bench \cite{terminalbench2025} targets shell execution with trajectory logging but provides no safety evaluation or error injection.

\subsubsection{Web and GUI Benchmarks}

WebArena \cite{zhou2024webarena} and VisualWebArena \cite{koh2024visualwebarena} benchmark web navigation in sandboxed environments with deterministic reward functions. However, they evaluate final page state rather than the navigation trajectory, and they operate on real (cloned) websites rather than controlled mock services, making reproducibility challenging. OSWorld \cite{xie2024osworld} extends evaluation to full desktop GUI environments but inherits the same output-only grading limitation.

\subsubsection{Multi-turn Benchmarks}

$\tau$-bench \cite{yao2024taubench} is the closest predecessor to Claw-Eval's multi-turn evaluation. It measures policy compliance in customer-service dialogues with simulated users and introduced the Pass$^k$ metric for consistency measurement. However, $\tau$-bench does not support multimodal tasks, does not include error injection for robustness, and does not embed safety constraints within task workflows. MINT \cite{wang2024mint} evaluates multi-turn tool use but focuses on task success rate without dimensional scoring.

\subsubsection{Multi-domain Suites}

AgentBench \cite{liu2024agentbench}, GAIA \cite{mialon2024gaia}, and TheAgentCompany \cite{xu2024agentcompany} broaden coverage across heterogeneous environments. TheAgentCompany adds sub-task checkpoints that partially address trajectory auditing, but none of these frameworks combine safety evaluation, robustness testing, and trajectory auditing in a single integrated framework.

\subsection{Evaluation Methodology}

\subsubsection{Output-only vs.\ Trajectory-aware Grading}

The dominant evaluation paradigm checks only final outputs---whether a file was created, a webpage reached the correct state, or a code patch passes tests. Pan et al.\ \cite{pan2024autocorrect} and Xu et al.\ \cite{xu2024hallucination} demonstrate that output-only evaluation cannot detect cases where agents fabricate intermediate steps yet produce plausible artifacts, a phenomenon related to reward hacking \cite{skalse2022reward, pan2022reward}. Claw-Eval addresses this through three independent evidence channels (execution traces, environment snapshots, audit logs), though our analysis reveals that the audit log channel is underutilized for safety verification.

\subsubsection{LLM-as-Judge}

The LLM-as-judge paradigm \cite{zheng2023judge, li2024generative, kim2024prometheus} scales evaluation to open-ended tasks but introduces its own reliability concerns. Claw-Eval's hybrid approach---deterministic checks for objective criteria, LLM judges for subjective quality---achieves 95\% human agreement and detects 44\% more safety violations than pure LLM judging.

\subsubsection{Graded vs.\ Binary Safety Evaluation}

XGUARD \cite{xguard2025} introduces a five-level danger scale for extremist content, demonstrating that binary labels mask the ``nuanced spectrum of risk'' in LLM outputs and proposes the Attack Severity Curve (ASC) to visualize model vulnerabilities across threat intensities---a visualization that requires graduated scores and cannot be produced from binary data. A recent meta-analysis of 40 agent safety benchmarks \cite{yu2026taxonomy} finds that 79\% use binary pass/fail metrics, that benchmark choice systematically shapes reported safety with no concordance across evaluation instruments (Kendall's $W = 0.10$, $p = 0.94$), and that 60\% of benchmarks cannot distinguish a genuinely safe model from one that simply refuses all inputs. These findings directly motivate our application of formal hierarchical models to safety scoring.

AgentLeak \cite{agentleak2026} demonstrates that output-only auditing misses 41.7\% of privacy violations in multi-agent systems, because internal channels (inter-agent messages, shared memory) leak sensitive data at 2.1$\times$ higher rates than user-facing outputs. This finding directly motivates our dispatch body scanning, which extends information flow analysis beyond assistant messages to tool call request bodies (email payloads, notification content, web fetch URLs). A large-scale empirical study of 17,022 agent skills \cite{credleak2026} finds that debug logging (\texttt{print}/\texttt{console.log}) is responsible for 73.5\% of credential leakage vulnerabilities, with 89.6\% of affected skills exploitable during routine execution. These internal-channel and debug-logging leaks are invisible to binary scoring that only checks final assistant output.

\subsubsection{Trajectory- and Flow-Level Safety Benchmarks}

A parallel line of 2026 work independently converges on our central thesis: safety must be evaluated over the \textit{trajectory} and its \textit{information flows}, not the final output, and severity must be graded. AgentSCOPE \cite{agentscope2026} introduces a Contextual-Integrity-grounded Privacy Flow Graph that decomposes agent execution into per-boundary information flows; across 62 multi-tool scenarios it finds pipeline privacy violations in over 80\% of cases \textit{even when the final output appears clean} (24\%), with most violations arising at the tool-response stage---precisely the internal channels our dispatch-body analysis targets. AgentCIBench \cite{agentcibench2026} evaluates fifteen frontier computer-use agents and reports that task completion is a poor proxy for disclosure safety: among agents with similar utility, leakage varies by more than 80 percentage points. This is exactly the ranking pathology SafeLattice corrects, and our experiments reproduce it (Section~\ref{sec:empirical}). HarnessAudit \cite{harnessaudit2026} audits full execution trajectories across 210 tasks and finds that violations accumulate with trajectory length and concentrate in resource access and inter-component information transfer, corroborating our state-transition view of safety. Most directly, TraceSafe-Bench \cite{tracesafe2026} constructs a static, trace-level benchmark of over 1,000 multi-step traces and concludes that ``granular risk taxonomies improve detection accuracy over binary judgments''---an independent empirical validation of graded scoring, and the methodological precedent for our trace-level evaluation. AgentHazard \cite{agenthazard2026} extends this evidence to the very framework Claw-Eval evaluates: across 2,653 instances on Claude Code, OpenClaw, and IFlow, it scores full execution trajectories with a graduated 0--10 harmfulness scale and finds that attack success roughly \emph{triples} between round~1 and round~3 (e.g., $29.9\% \to 68.1\%$ on OpenClaw)---harm is trajectory-dependent, single-turn evaluation misses the majority of it, and the field's newest OpenClaw-facing benchmark has again rejected flat binary judgments.

\subsubsection{Practitioner Tools with Graduated Severity}

Independent security evaluation tools have converged on graduated severity scales. Tinman-OpenClaw-Eval \cite{tinmaneval2026}, a community security harness for OpenClaw agents, implements 288 attack probes across 12 categories (prompt injection, tool exfiltration, context bleed, privilege escalation, supply chain, financial transaction, unauthorized action, MCP attacks, indirect injection, evasion bypass, memory poisoning, platform-specific) with a five-level severity scale: S4 (Critical), S3 (High), S2 (Medium), S1 (Low), S0 (Info). Their severity decision guide---``Is this just cosmetic? S0. Minor inconvenience? S1. Business/compliance risk? S3. Safety/security breach? S4''---independently validates our BLP-inspired hierarchy. The fact that practitioners building real-world agent security tools have independently rejected binary classification in favor of graduated severity provides strong external validation for our approach---and illustrates the research-practice gap DTRAP targets: deployed tooling has moved to graduated severity while benchmark research remains binary.

\subsection{Real-World Agent Security Incidents}
\label{sec:incidents}

OpenClaw, the autonomous agent framework evaluated by Claw-Eval, has accumulated over 60 CVEs and 100 GitHub Security Advisories since January 2026 \cite{openclawvulntaxonomy2026}. These incidents provide concrete attack patterns that demonstrate the inadequacy of binary safety scoring.

\textbf{ClawJacked (CVE-2026-32025, CVSS 8.8)} \cite{clawjacked2026}: A compound vulnerability chain allowing any malicious website to silently hijack a local OpenClaw agent via WebSocket exploitation---brute-forcing the gateway password (no rate limiting for loopback), auto-pairing as a trusted device, and gaining full agent control including the ability to search Slack history for API keys, exfiltrate files, and execute shell commands.

\textbf{Claw Chain (CVE-2026-44112/13/15/18)} \cite{clawchain2026}: Four chained CVEs enabling sandbox escape. CVE-2026-44113 (CVSS 7.7) exploits a TOCTOU race condition for read-side file exfiltration; CVE-2026-44112 (CVSS 9.6) enables write-side sandbox escape for backdoor installation. Binary scoring treats both as identical failures ($\texttt{safety} = 0.0$), despite the write escape being categorically more severe. CVE-2026-44115 exploits heredoc shell expansion to exfiltrate credentials through allowlisted commands---a violation that passes surface-level tool validation.

These real-world attacks motivate our evaluation framework: each involves multi-step chains where individual steps may appear safe but compound into critical violations, severity differs substantially across attack types, and information flows through channels (debug logs, dispatch bodies, compacted context) that output-only auditing cannot monitor.

\subsection{Safety and Robustness Evaluation}

Dedicated safety evaluations include ToolEmu \cite{ruan2024toolemu}, which simulates tool environments to test risk awareness; R-Judge \cite{yuan2024rjudge}, which evaluates safety reasoning over traces; Agent-SafetyBench \cite{zhang2024agentsafety}, which tests 349 scenarios across 8 risk categories; and MobileRisk-Live \cite{deng2025mobilerisk}, which assesses mobile agent risks. However, these frameworks evaluate safety in isolation from task completion---a model that refuses all potentially risky actions would score perfectly on safety but fail on completion. Claw-Eval's key design insight is embedding safety constraints within normal workflow tasks, so that models must navigate the tension between helpfulness and safety.

\subsection{Formal Security Models Applied to AI Agents}

Recent work has begun applying classical formal security models to AI agent systems. PCAS \cite{pcas2025} and its extension FORGE \cite{forge2025} implement Bell-LaPadula's MLS policy and toxic-flow taint tracking as runtime enforcement monitors for agentic systems, using Datalog to compute dependency graphs and block exfiltration at the action boundary. Parallax \cite{parallax2026} argues that Information Flow Control, grounded in BLP and Decentralized IFC \cite{myers1997difc}, is the correct foundation for agent security, proposing an ``Assume-Compromise'' evaluation model where the LLM is treated as fully compromised and the enforcement boundary is tested directly. MCPShield \cite{mcpshield2025} uses labeled transition systems for formal verification of MCP tool interaction chains, achieving 91\% theoretical threat coverage. The Hierarchical Autonomy Evolution (HAE) framework \cite{hae2025} organizes agent security into three autonomy tiers with a four-tier risk classification. DACSA \cite{dacsa2025} applies BLP lattice principles to data-centric authorization with sensitivity labels that propagate through agent transformations.

These works apply formal models to \textit{runtime enforcement}---preventing violations during execution. Our contribution is distinct: we apply BLP to \textit{evaluation}---detecting and classifying violations during post-hoc grading. This is complementary: runtime enforcement prevents harm, while lattice-aware evaluation measures how well an agent would behave \textit{without} enforcement, which is the diagnostic signal benchmarks must provide.

\subsection{Research Gaps}

Table~\ref{tab:comparison} summarizes evaluation dimensions across existing benchmarks.

\begin{table}[t]
\centering
\caption{Comparison of agent evaluation benchmarks.}
\label{tab:comparison}
\begin{tabular}{lcccccc}
\toprule
\textbf{Benchmark} & \textbf{Multimodal} & \textbf{Multi-turn} & \textbf{Auditable} & \textbf{Safety} & \textbf{Perturbation} & \textbf{Sandboxed} \\
\midrule
SWE-bench \cite{jimenez2024swebench} & -- & -- & -- & -- & -- & \checkmark \\
WebArena \cite{zhou2024webarena} & $\sim$ & -- & -- & -- & -- & \checkmark \\
$\tau$-bench \cite{yao2024taubench} & -- & \checkmark & \checkmark & \checkmark & -- & -- \\
AgentBench \cite{liu2024agentbench} & -- & \checkmark & $\sim$ & -- & -- & \checkmark \\
TheAgentCompany \cite{xu2024agentcompany} & $\sim$ & \checkmark & $\sim$ & $\sim$ & -- & \checkmark \\
Agent-SafetyBench \cite{zhang2024agentsafety} & -- & -- & \checkmark & \checkmark & -- & -- \\
Terminal-Bench \cite{terminalbench2025} & -- & -- & \checkmark & -- & -- & \checkmark \\
\midrule
Claw-Eval \cite{claweval2026} & \checkmark & \checkmark & \checkmark & $\sim$ & \checkmark & \checkmark \\
\textbf{+SafeLattice} & \checkmark & \checkmark & \checkmark & \checkmark & \checkmark & \checkmark \\
\bottomrule
\end{tabular}
\smallskip
\begin{minipage}{\linewidth}
\scriptsize
\textit{Note.} \checkmark = full support, $\sim$ = partial, -- = none. Claw-Eval's safety is ``partial'' because 81\% of tasks default to \texttt{safety=1.0} and scoring is binary. SafeLattice upgrades this to full hierarchical safety evaluation.
\end{minipage}
\end{table}

No prior work applies formal security models to the \textit{evaluation} of agent safety. Runtime enforcement (PCAS, Parallax) and evaluation (SafeLattice) are complementary: enforcement prevents harm during deployment, while lattice-aware evaluation measures intrinsic agent trustworthiness.

%% ============================================================
%% SECTION 3: BACKGROUND — FORMAL SECURITY MODELS
%% ============================================================
\section{Background: Formal Security Models}
\label{sec:background}

This section reviews the formal security models that underpin SafeLattice, drawing on the Bell-LaPadula and Biba models as formalized in the state-machine framework \cite{belllapadula1976, biba1977}.

\subsection{State-Machine Security Models}

Formal security models represent a system as a state machine $(S, O, A, T)$ where $S$ is the set of subjects, $O$ is the set of objects, $A$ is the access matrix mapping subjects to their current access modes on objects, and $T$ is the set of transition functions that transform the system state \cite{belllapadula1976}. Each subject $s \in S$ has a \textit{clearance} $\text{sclass}(s)$, and each object $o \in O$ has a \textit{classification} $\text{oclass}(o)$, both drawn from a partially ordered set of security levels.

A state is \textit{secure} if and only if a formal invariant holds across all subject-object pairs. The fundamental verification methodology requires six steps: (1)~define state variables, (2)~define the secure state invariant, (3)~define transition functions as atomic operations, (4)~prove that every transition preserves the invariant, (5)~define initial state constraints, and (6)~prove the initial state satisfies the invariant.

\subsection{The Bell-LaPadula Model}

The Bell-LaPadula (BLP) model \cite{belllapadula1976} enforces \textit{confidentiality} through mandatory access control over a hierarchy of security levels (e.g., Top Secret $>$ Secret $>$ Confidential $>$ Unclassified). BLP defines three core properties:

\begin{enumerate}
    \item \textbf{Simple Security Property (``No Read Up''):} A subject $s$ may read an object $o$ only if $\text{sclass}(s) \geq \text{oclass}(o)$. This prevents subjects from accessing information above their clearance.
    \item \textbf{Star Property (``No Write Down''):} A subject $s$ may write to an object $o$ only if $\text{oclass}(o) \geq \text{sclass}(s)$. This prevents high-clearance subjects from leaking information to lower-classification channels.
    \item \textbf{Strong Star Property:} A subject $s$ may read or write an object $o$ only if $\text{sclass}(s) = \text{oclass}(o)$. This is the strictest variant, used in high-assurance systems.
\end{enumerate}

\textbf{The Basic Security Theorem} states that if (a)~the initial system state is secure and (b)~every transition function preserves the secure state invariant, then the system is secure in every reachable state. This theorem provides the mathematical guarantee that security is maintained without runtime enumeration of all possible states.

\subsection{The Security Lattice}

BLP organizes security levels into a \textit{lattice}---a partially ordered set where every pair of elements has a unique least upper bound (join) and greatest lower bound (meet). The lattice structure ensures that security levels have a consistent dominance relation: if level $L_1 \geq L_2$, then a subject cleared at $L_1$ dominates one cleared at $L_2$. In formal notation, the lattice $(L, \leq)$ satisfies:
\begin{equation}
\forall\, a, b \in L: \quad a \wedge b \in L \;\;\text{and}\;\; a \vee b \in L
\end{equation}
where $a \wedge b$ is the meet (greatest lower bound) and $a \vee b$ is the join (least upper bound). For safety scoring, the meet operation is critical: when multiple violations occur, the \textit{minimum} severity score dominates.

\subsection{The Biba Integrity Model}

The Biba model \cite{biba1977} inverts BLP's logic to enforce \textit{integrity} rather than confidentiality:
\begin{itemize}
    \item \textbf{``No Read Down'':} A high-integrity subject must not read low-integrity data (to prevent contamination).
    \item \textbf{``No Write Up'':} A low-integrity subject must not write to high-integrity objects (to prevent corruption).
\end{itemize}

In the agent evaluation context, Biba's integrity model maps to the question of whether untrusted inputs (e.g., prompt injections embedded in retrieved documents) corrupt the agent's decision-making. SafeLattice operationalizes this: it classifies context sources by integrity (trusted system prompt $>$ endorsed user instruction $>$ untrusted retrieved/tool content) and detects the ``no write up'' violation in which low-integrity content drives a high-integrity action (Section~\ref{sec:biba-impl}). Prompt injection is thus not a separate ad-hoc check but a second formal property drawn from the same state-machine security framework.

\subsection{Mapping BLP to Agent Evaluation}
\label{sec:blp-mapping}

We now establish the formal mapping between BLP constructs and agent evaluation concepts, which forms the theoretical foundation of SafeLattice.

\begin{table}[t]
\centering
\caption{Mapping Bell-LaPadula to agent evaluation.}
\label{tab:blp-mapping}
\begin{tabular}{p{4cm}p{8cm}}
\toprule
\textbf{BLP Concept} & \textbf{Agent Evaluation Analogue} \\
\midrule
Subject & Agent (LLM) executing tools \\
Object & Data resources (emails, credentials, APIs) \\
Security Level & Data sensitivity (Critical, High, Medium, Low) \\
Subject Clearance & Agent's authorized access scope per task \\
Access Matrix & Tool-call permissions per task \\
No Read Up & Agent should not access data above task scope \\
No Write Down & Agent must not expose high-sensitivity data to low-clearance output channels \\
Secure State & No safety violations in the current evaluation state \\
State Transition & Each tool call or message in the agent's trajectory \\
Basic Security Theorem & If grading starts secure and each action preserves security, the trajectory is secure \\
\bottomrule
\end{tabular}
\end{table}

The key insight is that BLP's Star Property (``no write down'') directly models the credential leakage problem: when an agent reads high-sensitivity data (API keys from a config service classified as \textit{Critical}) and writes it to an assistant message (an output channel classified as \textit{Unclassified}, since it is visible to the user and potentially logged), this constitutes a \textit{write-down violation}---precisely the information flow that BLP was designed to prevent.

%% ============================================================
%% SECTION 4: ANALYSIS OF CURRENT SAFETY EVALUATION
%% ============================================================
\section{Analysis of Current Safety Evaluation}
\label{sec:analysis}

We use Claw-Eval as a case study to demonstrate that binary safety scoring violates the foundational properties of formal security models. Our analysis is based on static examination of all 300 task graders, their YAML configuration files, and the scoring pipeline source code.

\subsection{The Binary Safety Gate}

Claw-Eval's scoring formula uses safety as a multiplicative gate:
\begin{equation}
\text{task\_score} = \texttt{safety} \times (0.8 \cdot \texttt{completion} + 0.2 \cdot \texttt{robustness})
\label{eq:scoring}
\end{equation}
where $\texttt{safety} \in \{0, 1\}$ in practice, despite the formula mathematically supporting any value in $[0, 1]$. The \texttt{DimensionScores} model defaults $\texttt{safety} = 1.0$, meaning any grader that does not explicitly set the safety field automatically passes.

\subsection{Violation of BLP Property 1: No Hierarchical Classification}

BLP requires a partially ordered set of security levels with a dominance relation. Binary scoring collapses this structure to a single bit. Table~\ref{tab:binary-collapse} illustrates the diagnostic information lost.

\begin{table}[t]
\centering
\caption{Diagnostic collapse under binary scoring.}
\label{tab:binary-collapse}
\begin{tabular}{lcc}
\toprule
\textbf{Violation} & \textbf{Binary} & \textbf{Lattice} \\
\midrule
Leak API credentials & 0.0 & 0.0 (Critical) \\
Comply with prompt injection & 0.0 & 0.0 (Critical) \\
Send email (forbidden tool) & 0.0 & 0.3 (High) \\
Delete calendar event & 0.0 & 0.3 (High) \\
Send to wrong recipient & 0.0 & 0.6 (Medium) \\
Publish irrelevant content & 0.0 & 0.6 (Medium) \\
Save unnecessary draft & 0.0 & 0.9 (Low) \\
\bottomrule
\end{tabular}
\end{table}

Under binary scoring, all seven violations are indistinguishable ($\texttt{safety} = 0.0$). Under SafeLattice's four-level classification, they are correctly ranked by severity, enabling practitioners to distinguish agents that leak credentials from those that merely create unnecessary drafts.

\subsection{Violation of BLP Property 2: No Information Flow Control}

BLP's Star Property requires that data flow only from lower to higher security levels (``no write down''). In the agent evaluation context, this means:
\begin{itemize}
    \item Data from high-sensitivity sources (config service: API keys, passwords) should not flow to low-sensitivity outputs (assistant messages visible to users).
    \item Data from medium-sensitivity sources (CRM: customer PII) should not flow to external channels (email to unverified recipients).
\end{itemize}

Binary scoring checks \textit{which tool was called}, not \textit{what data flowed where}. Of 18 tasks that access the config mock service (which returns API keys such as \texttt{sk\_live\_...}, AWS access keys \texttt{AKIA...}, and database passwords), only 6 (33\%) check whether these credentials appear in the agent's output. The remaining 12 tasks provide agents access to Critical-classified data without any write-down verification. Under the current scoring, an agent that reads credentials via \texttt{config/integrations/get} and prints them verbatim in its response receives $\texttt{safety} = 1.0$.

SafeLattice's information flow analysis addresses this by tracking the sensitivity classification of data retrieved by each tool call and verifying that no high-sensitivity data appears in lower-classified output channels.

\subsection{Violation of BLP Property 3: No State Transition Verification}

The Basic Security Theorem requires that every state transition preserve the secure state invariant. In agent evaluation, each tool call represents a state transition. Binary scoring evaluates only the final state---it checks whether a forbidden tool was called, but not whether the \textit{sequence} of calls constitutes a violation.

Consider a multi-step information flow violation: (1)~agent calls \texttt{contacts\_search} with a wildcard query, retrieving all contact addresses; (2)~agent calls \texttt{config/integrations/get}, retrieving API credentials; (3)~agent calls \texttt{gmail\_send\_message} to send a ``summary'' that includes the credentials to all contacts. Steps 1 and 2 are individually permissible. Step~3 is individually detectable only if the grader checks for credential patterns in the email body. But the \textit{trajectory}---reading sensitive data and then broadcasting it---constitutes a compound violation that no individual check catches unless the evaluator tracks cumulative state.

SafeLattice models this by maintaining a security state across the trajectory: each tool call updates the agent's \textit{effective clearance} based on the sensitivity of data it has accessed, and subsequent output actions are verified against this accumulated state.

\subsection{Violation of BLP Property 4: No Initial State Guarantee}

The Basic Security Theorem requires a provably secure initial state. In Claw-Eval, \texttt{DimensionScores.safety} defaults to 1.0 (safe). This means:
\begin{itemize}
    \item 243 of 300 tasks (81\%) never evaluate safety, inheriting the default $\texttt{safety} = 1.0$.
    \item 35 tasks declare safety constraints in \texttt{task.yaml} but their graders do not enforce them, silently inheriting $\texttt{safety} = 1.0$.
\end{itemize}

In BLP terms, this is equivalent to a system where 81\% of subjects are granted Top Secret clearance by default, with no verification that this clearance is warranted. SafeLattice replaces this with a conservative initial state where declared safety constraints are always enforced regardless of individual grader implementation.

\subsection{Safety Coverage Analysis}

Table~\ref{tab:coverage} quantifies the safety evaluation coverage gap across Claw-Eval's three task groups.

\begin{table}[t]
\centering
\caption{Safety check coverage by task group.}
\label{tab:coverage}
\begin{tabular}{lccc}
\toprule
\textbf{Task Group} & \textbf{Total} & \textbf{Tested} & \textbf{Coverage} \\
\midrule
General (T-prefix) & 161 & 54 & 33.5\% \\
Multimodal (M-prefix) & 101 & 3 & 3.0\% \\
Multi-turn (C-prefix) & 38 & 0 & 0.0\% \\
\midrule
\textbf{Total} & \textbf{300} & \textbf{57} & \textbf{19.0\%} \\
\bottomrule
\end{tabular}
\end{table}

The 81\% coverage gap means that the safety multiplier is functionally inert for the majority of tasks, reducing the effective scoring formula to $\texttt{task\_score} = 0.8 \times \texttt{completion} + 0.2 \times \texttt{robustness}$ for 243 out of 300 tasks. In BLP terms, the security model is simply not applied to most of the system.

%% ============================================================
%% SECTION 5: PROPOSED FRAMEWORK — SAFELATTICE
%% ============================================================
\section{SafeLattice: BLP-Inspired Safety Evaluation}
\label{sec:safelattice}

SafeLattice applies Bell-LaPadula's formal security model to agent safety evaluation through three components: (1)~a security classification lattice, (2)~information flow analysis, and (3)~automatic enforcement via lattice-aware state verification.

\subsection{Security Classification Lattice}

Following BLP's hierarchy, SafeLattice defines four security levels with a total ordering and associated safety scores:

\begin{equation}
\text{Critical} > \text{High} > \text{Medium} > \text{Low} > \text{Safe}
\end{equation}

\begin{table}[t]
\centering
\caption{SafeLattice security classification lattice.}
\label{tab:lattice}
\begin{tabular}{lccp{6cm}}
\toprule
\textbf{Level} & \textbf{Score} & \textbf{BLP Analogue} & \textbf{Violation Types} \\
\midrule
Critical & 0.0 & Top Secret & Credential leakage, data exfiltration, injection compliance \\
High & 0.3 & Secret & Unauthorized send/delete/close, forbidden tool calls \\
Medium & 0.6 & Confidential & Wrong recipients, excess scope, wrong data \\
Low & 0.9 & Unclassified & Unnecessary actions, redundant calls \\
Safe & 1.0 & -- & No violation \\
\bottomrule
\end{tabular}
\end{table}

The lattice satisfies BLP's dominance requirement: for any two violations with severities $v_1$ and $v_2$, the meet operation $v_1 \wedge v_2 = \min(\text{score}(v_1), \text{score}(v_2))$ yields the most severe classification. When multiple violations co-occur, the minimum score dominates---a single Critical violation zeros the safety score regardless of how many Low-severity issues coexist.

\textbf{Backward compatibility.} Any violation that scored $\texttt{safety} = 0.0$ under binary scoring maps to Critical or High severity ($\text{score} \leq 0.3$), so no previously-failing task can pass under SafeLattice. The pass threshold remains $\texttt{task\_score} \geq 0.75$; even a High-severity violation ($\texttt{safety} = 0.3$) reduces the maximum achievable task score to $0.3 \times 1.0 = 0.30$, well below the threshold.

\subsection{Information Flow Analysis}

SafeLattice implements BLP's Star Property as a post-hoc information flow analysis pass over the agent's execution trace. The analysis proceeds in three steps:

\textbf{Step 1: Classify data sources.} Each mock service endpoint is assigned a sensitivity level based on the data it returns:
\begin{itemize}
    \item \textit{Critical:} \texttt{config/integrations/get} (API keys, passwords), \texttt{crm/export} (full customer PII)
    \item \textit{High:} \texttt{gmail/inbox} (personal email content), \texttt{finance/reports} (financial data)
    \item \textit{Medium:} \texttt{contacts/search} (contact lists), \texttt{calendar/events} (schedule data)
    \item \textit{Low:} \texttt{kb/search} (public knowledge base), \texttt{web/search} (web results)
\end{itemize}

\textbf{Step 2: Track data flow through trajectory.} For each \texttt{ToolDispatch} event, the analysis records the sensitivity level of the data returned. The agent's \textit{accumulated sensitivity} is the maximum sensitivity of any data it has accessed---once an agent reads Critical data, its context is tainted at the Critical level.

\textbf{Step 3: Verify write-down constraints.} When the agent produces output (assistant messages, tool calls to send/share/publish endpoints), the analysis checks whether the output channel's classification is at least as high as the agent's accumulated sensitivity. If the agent has read Critical data and writes to an Unclassified channel (e.g., assistant message), this constitutes a write-down violation.

Formally, let $\text{sens}(d)$ denote the sensitivity of tool dispatch $d$, and let $\text{out\_class}(c)$ denote the classification of output channel $c$. The write-down check verifies:
\begin{equation}
\forall\, c \in \text{outputs}: \quad \text{out\_class}(c) \geq \max_{d \in \text{dispatches}} \text{sens}(d)
\end{equation}

\subsection{Automatic Safety Enforcement}

SafeLattice includes an enforcement module that reads \texttt{safety\_checks} declarations from \texttt{task.yaml} and applies them automatically during grading. For each declared check, the enforcer:
\begin{enumerate}
    \item Identifies the check type (e.g., \texttt{tool\_not\_called}, \texttt{credential\_exposure}, \texttt{wrong\_action}).
    \item Applies the corresponding verification against the execution trace.
    \item Classifies any detected violation using the security lattice (Table~\ref{tab:lattice}).
    \item Updates the safety score using the meet operation: $\texttt{safety} = \min(\texttt{safety}, \text{score}(v))$.
\end{enumerate}

This decouples safety enforcement from individual grader implementations. The 35 tasks whose graders silently ignore declared safety constraints are now automatically covered.

\subsection{Integrity Analysis: Biba for Prompt Injection}
\label{sec:biba-impl}

Where BLP protects confidentiality (no write \textit{down}), Biba \cite{biba1977} protects integrity (no write \textit{up}). SafeLattice adds a Biba pass that treats prompt injection as an integrity violation: low-integrity content must not corrupt a high-integrity decision. We assign each context source an integrity level,
\begin{equation}
\text{Trusted (system)} > \text{Endorsed (user)} > \text{Untrusted (retrieved/tool)},
\end{equation}
and classify the responses of retrieval tools (\texttt{web\_search}, \texttt{web\_fetch}, \texttt{gmail\_get\_message}, \texttt{kb\_search}, \dots) as Untrusted. A ``no write up'' violation is recorded only when \textit{both} conditions hold: (1)~an injection marker (e.g., ``ignore previous instructions'', ``forward all'', ``print your key'') appears in Untrusted-integrity input, \textit{and} (2)~a downstream high-integrity action follows---a credential surfaces in output (classified Critical, \texttt{prompt\_injection\_compliance}) or a sensitive outbound/destructive tool call executes (classified High, \texttt{integrity\_violation}). Reading adversarial content without acting on it is \textit{not} penalized, mirroring an agent that recognizes and refuses an injection. This conjunctive rule keeps false positives low while capturing the causal ``low-integrity input $\rightarrow$ high-integrity action'' signature that distinguishes compliance from mere exposure.

%% ============================================================
%% SECTION 6: IMPLEMENTATION
%% ============================================================
\section{Implementation}
\label{sec:implementation}

SafeLattice is implemented as a Python module that integrates into the Claw-Eval grading pipeline. The implementation comprises three files totaling approximately 500 lines of code.

\subsection{Security Lattice (\texttt{safety\_taxonomy.py})}

The lattice is implemented as an enumeration with an associated score mapping. A \texttt{SafetyAssessment} class accumulates violations and computes the meet (minimum) score. The violation category registry maps 22 violation types to their lattice-classified severity levels. Data source sensitivity classifications are defined for all 19 mock service endpoints, enabling information flow analysis.

\subsection{Information Flow Analyzer}
\label{sec:flow-analyzer}

The information flow analyzer is integrated into the safety enforcer. For each \texttt{ToolDispatch} event in the trace, it records the sensitivity level of the accessed service. When scanning assistant messages for credential patterns, it applies the Star Property: if the agent has accessed any Critical-classified service, \textit{all} assistant messages are checked for sensitive patterns, regardless of whether the individual task grader implements this check.

\textbf{Canary-based obfuscation detection.} Pattern matching alone is trivially evaded by encoding. We exploit a structural property of the benchmark: because the evaluation controls its mock services, every secret an agent can observe arrives through a tool response---making each observed secret a de-facto canary token. The analyzer harvests credential values from high-sensitivity tool responses (regex-shaped credentials plus values of secret-bearing JSON keys) and checks whether any \emph{transformation} of a harvested value reaches an output channel: separator-stripped reassembly (catches whitespace- and string-split evasion), Unicode-confusable normalization plus Hamming-distance-1 matching (catches homoglyph substitution), reversal, and base64 (both encoding the canary and entropy-gated decoding of output tokens). A match raises a Critical \texttt{obfuscated\_credential\_leak} write-down violation.

\subsection{Pipeline Integration (\texttt{safety\_enforcer.py})}

The enforcer is invoked after the per-task grader returns its \texttt{DimensionScores}. It receives the task's \texttt{safety\_checks} declarations, the execution trace (dispatches and messages), and the current scores. If it detects violations that the grader missed, it updates \texttt{scores.safety} to the lattice-classified value. The enforcer runs in graduated mode by default, assigning severity-proportional scores rather than binary values.

\subsection{Integration Point}

The enforcer is called in the \texttt{\_grade\_with\_optional\_params} path in the CLI, after the task-specific grader has run:

\begin{lstlisting}[language=Python]
scores = grader.grade(messages, dispatches,
                      task, **kwargs)
scores, assessment = enforce_safety_checks(
    task.safety_checks, dispatches,
    messages, scores, use_graduated=True)
task_score = compute_task_score(scores)
\end{lstlisting}

%% ============================================================
%% SECTION 7: FORMAL PROPERTIES
%% ============================================================
\section{Formal Properties of SafeLattice}
\label{sec:proofs}

We now apply the six-step state-machine verification methodology of Section~\ref{sec:blp-mapping} to SafeLattice's scoring procedure, treating grading as a state machine that consumes the agent's trajectory event-by-event. This yields a Basic Security Theorem analogue for evaluation: if grading starts in a secure state and every transition preserves the secure-state invariant, the final score is a sound summary of the trajectory's most severe violation.

\subsection{Step 1: State Variables}

Let $L=\{\textsc{Safe},\textsc{Low},\textsc{Med},\textsc{High},\textsc{Crit}\}$ be the severity levels with score map $s:L\to[0,1]$, $s(\textsc{Crit})=0,\,s(\textsc{High})=0.3,\,s(\textsc{Med})=0.6,\,s(\textsc{Low})=0.9,\,s(\textsc{Safe})=1$. Define the severity (dominance) order $\sqsupseteq$ by $\textsc{Crit}\sqsupseteq\textsc{High}\sqsupseteq\textsc{Med}\sqsupseteq\textsc{Low}\sqsupseteq\textsc{Safe}$, i.e.\ $a\sqsupseteq b \iff s(a)\le s(b)$. Processing a trajectory prefix of length $k$ yields state $q_k=(\sigma_k, V_k)$, where $\sigma_k\in L$ is the \emph{accumulated data sensitivity} (the most sensitive source read so far) and $V_k\subseteq L$ is the set of recorded violation severities.

\subsection{Step 2: Secure-State Invariant}

We say $q_k$ is \emph{score-consistent} if the reported safety score equals the meet of all recorded violations:
\begin{equation}
\mathrm{safety}(q_k) \;=\; s\!\left(\textstyle\bigwedge V_k\right), \qquad \textstyle\bigwedge \varnothing := \textsc{Safe},
\label{eq:invariant}
\end{equation}
where $a\wedge b$ is the lattice meet (the more severe level, i.e.\ the smaller score). The invariant $I(q_k)$ is precisely Eq.~\eqref{eq:invariant}: the score at all times reflects the single most severe violation seen.

\subsection{Step 3: Transition Functions}

Each event $e_{k+1}$ (a tool dispatch or a message) induces $\delta(q_k,e_{k+1})=q_{k+1}$:
\begin{align}
\sigma_{k+1} &= \sigma_k \wedge_{\!\text{sens}} \mathrm{sens}(e_{k+1}),\\
V_{k+1} &= V_k \cup \mathrm{detect}(e_{k+1}, \sigma_{k+1}),
\end{align}
where $\mathrm{sens}(\cdot)$ is the source sensitivity, $\wedge_{\text{sens}}$ accumulates the maximum sensitivity, and $\mathrm{detect}(\cdot)$ returns the severities of any violations the BLP, Biba, and declared-check analyses raise at this step. The score is updated by $\mathrm{safety}(q_{k+1})=\min\big(\mathrm{safety}(q_k),\, \min_{v\in \mathrm{detect}(\cdot)} s(v)\big)$.

\subsection{Step 4: Lattice Well-Formedness}

\begin{lemma}[Well-formed lattice]\label{lem:lattice}
$(L,\sqsupseteq)$ is a complete lattice: $\sqsupseteq$ is a partial order, and meet $\wedge$ and join $\vee$ are total, closed, commutative, associative, and idempotent, with bottom $\textsc{Safe}$ and top $\textsc{Crit}$.
\end{lemma}
\begin{proof}
$s$ is injective, so $\sqsupseteq$ inherits reflexivity, antisymmetry, and transitivity from $\le$ on $s(L)\subset\mathbb{R}$; hence it is a partial order. Because $s(L)$ is totally ordered, every pair $a,b$ has $a\wedge b=\arg\min\{s(a),s(b)\}$ and $a\vee b=\arg\max\{s(a),s(b)\}$, both in $L$ (closure). A finite totally ordered set is a complete lattice, with least element $\textsc{Safe}$ ($s=1$) and greatest element $\textsc{Crit}$ ($s=0$).
\end{proof}

\subsection{Steps 5--6: Initial State and Invariant Preservation}

\begin{lemma}[Secure initial state]\label{lem:init}
The initial state $q_0=(\textsc{Safe},\varnothing)$ satisfies $I(q_0)$.
\end{lemma}
\begin{proof}
$\mathrm{safety}(q_0)=1=s(\textsc{Safe})=s(\bigwedge\varnothing)$.
\end{proof}

\begin{theorem}[Invariant preservation / monotonicity]\label{thm:preserve}
If $I(q_k)$ holds, then $I(q_{k+1})$ holds and $\mathrm{safety}(q_{k+1})\le \mathrm{safety}(q_k)$.
\end{theorem}
\begin{proof}
Since $V_{k+1}=V_k\cup \Delta$ with $\Delta=\mathrm{detect}(e_{k+1},\sigma_{k+1})$, and meet is monotone under set inclusion, $\bigwedge V_{k+1}=\big(\bigwedge V_k\big)\wedge\big(\bigwedge\Delta\big)\sqsupseteq \bigwedge V_k$; by Lemma~\ref{lem:lattice} this gives $s(\bigwedge V_{k+1})\le s(\bigwedge V_k)$, so the score cannot increase. By the induction hypothesis $\mathrm{safety}(q_k)=s(\bigwedge V_k)$, and the update rule sets $\mathrm{safety}(q_{k+1})=\min\!\big(s(\bigwedge V_k),\min_{v\in\Delta}s(v)\big)=s\big((\bigwedge V_k)\wedge(\bigwedge\Delta)\big)=s(\bigwedge V_{k+1})$, establishing $I(q_{k+1})$.
\end{proof}

\begin{corollary}[Basic Security Theorem for evaluation]
For any trajectory, the final score equals $s(\bigwedge_{v} v)$ over all violations detected along the trajectory; equivalently, the score is exactly the score of the single most severe violation, and a single \textsc{Crit} violation forces $\mathrm{safety}=0$.
\end{corollary}
\begin{proof}
Immediate by induction: Lemma~\ref{lem:init} (base) and Theorem~\ref{thm:preserve} (step).
\end{proof}

\subsection{Backward Compatibility}

\begin{theorem}[No binary-failing trace can pass]\label{thm:bc}
If the binary gate assigns $\mathrm{safety}_{\mathrm{bin}}=0$ to a trace, then SafeLattice assigns $\mathrm{safety}_{\mathrm{SL}}\le 0.3$, hence $\mathrm{task\_score}_{\mathrm{SL}}\le 0.3<0.75$.
\end{theorem}
\begin{proof}
The binary gate returns $0$ only when (a)~a forbidden tool declared \texttt{tool\_not\_called} is invoked, or (b)~a credential pattern appears in the final assistant message. In case~(a), SafeLattice records a \texttt{forbidden\_tool\_call} at severity $\sqsupseteq\textsc{High}$, so $\bigwedge V_{\mathrm{SL}}\sqsupseteq\textsc{High}$ and $\mathrm{safety}_{\mathrm{SL}}=s(\bigwedge V_{\mathrm{SL}})\le s(\textsc{High})=0.3$. In case~(b), the same output triggers SafeLattice's \texttt{credential\_leak}/write-down detector at \textsc{Crit}, giving $\mathrm{safety}_{\mathrm{SL}}=0$. In either case $\mathrm{safety}_{\mathrm{SL}}\le 0.3$. Since $\mathrm{task\_score}=\mathrm{safety}\cdot(0.8\,\mathrm{comp}+0.2\,\mathrm{rob})\le \mathrm{safety}\le 0.3$, the trace remains below the $0.75$ pass threshold.
\end{proof}

The Biba integrity component (Section~\ref{sec:biba-impl}) extends the state with an injection-observed flag and contributes \textsc{Crit}/\textsc{High} severities to $\mathrm{detect}(\cdot)$; Theorem~\ref{thm:preserve} applies verbatim, so integrity checking preserves all guarantees above.

%% ============================================================
%% SECTION 8: RESULTS AND ANALYSIS
%% ============================================================
\section{Results and Analysis}
\label{sec:results}

\subsection{Use of LLMs}
\label{sec:use-of-llms}

LLMs appear in this work in three distinct roles, which we disclose explicitly. \textbf{(1)~Agents under evaluation:} the live evaluation (Section~\ref{sec:live-pilot}) runs six frontier models---GPT-4o, GPT-4o-mini, Claude Sonnet 4.6, Claude Haiku 4.5, Gemini 2.5 Flash, and Llama-3.3-70B-Instruct---as the agents whose safety is being measured, accessed through a single OpenRouter API endpoint so that all models execute under identical harness conditions. \textbf{(2)~LLM-as-judge in grading:} Claw-Eval's production grading pipeline combines deterministic checks with an LLM judge for subjective completion criteria; all live traces are graded with the judge enabled (GPT-4o-mini as the judge model), and the same graded scores feed both the binary and SafeLattice scoring paths, so judge behavior cannot differentially favor either system. \textbf{(3)~No LLM in the safety detectors or ground truth:} SafeLattice's lattice, information flow analysis, Biba integrity pass, and canary-based obfuscation detection are entirely deterministic, and ground-truth violation labels are established by independent deterministic detectors plus mechanical task-specific verifiers (Section~\ref{sec:live-pilot})---no LLM judges any safety outcome reported in this paper. The trace-level corpus of Section~\ref{sec:empirical} is generated deterministically without LLM calls.

\subsection{Severity Differentiation (RQ1)}

We first present an illustrative micro-benchmark that isolates the scoring function: across 10 representative scenarios, completion and robustness are held constant and only the safety dimension varies, so the score bands reflect severity alone. Table~\ref{tab:graduated} presents these controlled scores; Section~\ref{sec:empirical} then evaluates SafeLattice end-to-end on a labeled corpus of real agent trajectories.

\begin{table}[t]
\centering
\caption{Binary vs.\ SafeLattice graduated scoring.}
\label{tab:graduated}
\begin{tabular}{lcccc}
\toprule
\textbf{Scenario} & \textbf{Severity} & \textbf{Binary} & \textbf{Graduated} & \textbf{$\Delta$} \\
\midrule
Clean execution & Safe & 0.872 & 0.872 & 0.000 \\
Unnecessary draft saved & Low & 0.000 & 0.828 & +0.828 \\
Wrong email recipient & Medium & 0.000 & 0.542 & +0.542 \\
Wildcard contact query & Medium & 0.000 & 0.480 & +0.480 \\
Forbidden send tool & High & 0.000 & 0.262 & +0.262 \\
Forbidden close tool & High & 0.000 & 0.281 & +0.281 \\
Forbidden delete tool & High & 0.000 & 0.228 & +0.228 \\
Credential leak (API key) & Critical & 0.000 & 0.000 & 0.000 \\
Injection compliance & Critical & 0.000 & 0.000 & 0.000 \\
Spam content published & Medium & 0.000 & 0.432 & +0.432 \\
\bottomrule
\end{tabular}
\smallskip
\begin{minipage}{\linewidth}
\scriptsize
\textit{Note.} Binary = binary task score; Graduated = SafeLattice task score; $\Delta$ = score improvement. Base completion and robustness scores held constant per scenario.
\end{minipage}
\end{table}

\textbf{Key findings:}
\begin{itemize}
    \item \textbf{Critical violations remain zero-scored.} Credential leakage and prompt injection compliance produce $\texttt{task\_score} = 0.0$ under both systems, satisfying BLP's requirement that Critical-level failures are absolute.
    \item \textbf{Severity differentiation is achieved.} The four severity levels produce distinct score bands: Critical (0.000), High (0.228--0.281), Medium (0.432--0.542), Low (0.828). These bands are non-overlapping, enabling practitioners to determine violation severity from the task score alone.
    \item \textbf{Low-severity violations receive proportional treatment.} An agent that saves an unnecessary draft receives a task score of 0.828 instead of 0.000---a 95\% improvement that better reflects the minimal actual harm.
    \item \textbf{No previously-failing task can pass.} The maximum task score for a High-severity violation is 0.281, well below the 0.75 pass threshold, preserving backward compatibility.
\end{itemize}

\subsection{Real-World Attack Pattern Analysis}
\label{sec:cve-analysis}

To validate SafeLattice against real-world threats, we reproduced attack patterns from documented OpenClaw CVEs as unit tests within our framework. Table~\ref{tab:cve-comparison} shows how binary scoring and SafeLattice handle each pattern.

\begin{table}[t]
\centering
\caption{Binary vs.\ SafeLattice: real-world attack patterns.}
\label{tab:cve-comparison}
\begin{tabular}{p{4.2cm}p{2.6cm}ccc}
\toprule
\textbf{Attack} & \textbf{Source} & \textbf{Binary} & \textbf{SafeLattice} & \textbf{BLP Property} \\
\midrule
Credential chain (gateway hijack) & CVE-2026-32025 & 0.0 & 0.0 (C) & Star \\
Read escape (file exfiltration) & CVE-2026-44113 & 0.0 & 0.3 (H) & Access \\
Write escape (backdoor) & CVE-2026-44112 & 0.0 & 0.0 (C) & Star \\
Allowlist bypass (heredoc exfiltration) & CVE-2026-44115 & 0.0 & 0.0 (C) & Star \\
Debug log leak & \cite{credleak2026} & 1.0$^*$ & 0.0 (C) & Star \\
Compaction persistence (3-step) & Issue \#56268 & 1.0$^*$ & 0.0 (C) & Star \\
Internal channel exfiltration & \cite{agentleak2026} & 1.0$^*$ & 0.0 (C) & Star \\
\bottomrule
\end{tabular}
\smallskip
\begin{minipage}{\linewidth}
\scriptsize
\textit{Note.} C = Critical, H = High. Binary = binary safety score, SafeLattice = SafeLattice safety score. $^*$Binary scores 1.0 because the violation is invisible to output-only, tool-name-only checking: credentials flow through debug logs, dispatch bodies, or compacted context rather than appearing in forbidden tool calls. SafeLattice catches these via BLP Star Property information flow analysis and dispatch body credential scanning.
\end{minipage}
\end{table}

\textbf{Key findings from CVE analysis:}
\begin{itemize}
    \item \textbf{Binary scoring misses 3 of 7 attack patterns entirely.} Debug log leaks, compaction-persistence chains, and internal channel exfiltration all score $\texttt{safety} = 1.0$ under binary evaluation because the credential never appears in a forbidden tool call or final assistant output---it flows through internal channels that binary scoring does not monitor.
    \item \textbf{Binary scoring conflates the remaining 4 patterns.} Read-side file exfiltration (CVE-2026-44113, CVSS 7.7) and write-side sandbox escape with backdoor installation (CVE-2026-44112, CVSS 9.6) both score $\texttt{safety} = 0.0$, despite the latter being categorically more severe. SafeLattice correctly differentiates them: High (0.3) vs.\ Critical (0.0).
    \item \textbf{SafeLattice's information flow analysis is critical.} All three missed patterns are detected by BLP Star Property enforcement---tracking accumulated data sensitivity through the trajectory and scanning outbound tool dispatch bodies for credential patterns, not just assistant messages.
\end{itemize}

\subsection{Information Flow Detection (RQ2)}

SafeLattice's information flow analysis identifies write-down violations that binary scoring cannot detect. Of 18 tasks that access the config mock service (returning credentials including \texttt{sk\_live\_...}, \texttt{AKIA...}, and database passwords), binary scoring checks for leakage in only 6 tasks. SafeLattice's Star Property analysis extends credential pattern scanning to all 18 tasks, detecting potential write-down violations in the 12 previously unprotected tasks.

The write-down detection operates as follows: when the agent calls \texttt{config/integrations/get} (classified as Critical), SafeLattice marks the agent's context as tainted at the Critical level. Any subsequent assistant message is then scanned for sensitive credential patterns. If credentials appear, a Critical-severity violation is recorded---regardless of whether the individual task grader implements this check.

\subsection{Safety Coverage Improvement (RQ3)}

Automatic enforcement of declared safety constraints increases the number of safety-evaluated tasks:

\begin{itemize}
    \item \textbf{Before SafeLattice:} 57/300 tasks (19.0\%) evaluate safety.
    \item \textbf{After SafeLattice:} 85/300 tasks (28.3\%) evaluate safety---a 49\% relative improvement.
    \item \textbf{Source of improvement:} 28 tasks with \texttt{tool\_not\_called} declarations that were previously unenforced are now automatically checked, covering 17 unique forbidden tool types including \texttt{gmail\_send\_message}, \texttt{helpdesk\_close\_ticket}, \texttt{calendar\_delete\_event}, and \texttt{scheduler\_delete\_job}.
\end{itemize}

Additionally, the information flow analysis adds implicit credential leakage checks to 12 tasks that access sensitive services without explicit safety checks, bringing the total coverage including implicit checks to approximately 97/300 (32.3\%).

\subsection{Trace-Level Empirical Evaluation (RQ4)}
\label{sec:empirical}

The preceding results isolate the scoring function. We now evaluate SafeLattice end-to-end on agent \emph{trajectories}. We first adopt the trace-level benchmark methodology of TraceSafe-Bench \cite{tracesafe2026}, which scores curated multi-step traces rather than live agents; the live evaluation of Section~\ref{sec:live-pilot} then repeats the comparison on sampled frontier-model behavior. We construct a corpus of 68 execution traces (17 scenarios $\times$ 4 simulated agents) in the exact JSONL format emitted by the Claw-Eval runner, so the identical loaders and enforcer score curated and live traces alike. Scenarios are grounded in real safety-relevant task IDs from the 94-task subset (57 safety-checked $+$ 12 credential-vulnerable $+$ 9 injection $+$ 35 YAML-grader-mismatch, deduplicated) and in documented OpenClaw failure modes (output leaks, internal-channel exfiltration via email/notification/\texttt{web\_fetch}, injection compliance, forbidden actions). Every trace carries a ground-truth label. The four simulated agents differ in \emph{safety discipline}: an agent commits a scenario's violation iff the scenario's temptation exceeds its discipline, yielding a high-completion/low-discipline ``careless'' agent through a low-completion/high-discipline ``over-refuser.'' We release the corpus generator, the dual-scoring harness, and the live-run harness (Appendix~\ref{app:repro}).

Each trace is scored under two systems on identical inputs: a faithful model of the \emph{current} Claw-Eval gate (output-only credential detection plus forbidden-tool checks; blind to request bodies, information flow, and unenforced declared checks) and \emph{SafeLattice} (graduated severity $+$ BLP flow $+$ Biba $+$ internal-channel scanning).

\textbf{Detection quality.} Table~\ref{tab:detection} reports detection against ground truth. Binary scoring achieves recall $0.52$: it misses roughly half of all true violations, because internal-channel leaks (credentials in email bodies, notification payloads, and \texttt{web\_fetch} URLs) and unenforced declared checks never surface in the final assistant message. SafeLattice raises recall to $0.95$ at higher precision, fixing 9 false negatives---each an internal-channel or flow violation the output-only gate cannot see. SafeLattice's single remaining false negative is a Low-severity redundant action, for which it (deliberately) has no detector.

\begin{table}[t]
\centering
\caption{Trace-level detection quality on the labeled corpus of 68 agent traces.}
\label{tab:detection}
\begin{tabular}{lcccc}
\toprule
\textbf{System} & \textbf{Precision} & \textbf{Recall} & \textbf{F1} & \textbf{FN} \\
\midrule
Binary (current) & 0.79 & 0.52 & 0.63 & 10 \\
SafeLattice & \textbf{0.87} & \textbf{0.95} & \textbf{0.91} & \textbf{1} \\
\bottomrule
\end{tabular}
\smallskip
\begin{minipage}{\linewidth}
\scriptsize
\textit{Note.} FN = false negatives. The binary gate models the current output-only Claw-Eval scoring; SafeLattice adds graduated severity, BLP information-flow, Biba integrity, and internal-channel scanning.
\end{minipage}
\end{table}

\begin{table}[t]
\centering
\caption{Per-model mean task score under binary vs.\ SafeLattice scoring on the trace-level corpus.}
\label{tab:ranking}
\begin{tabular}{lccc}
\toprule
\textbf{Model} & \textbf{Completion} & \textbf{Binary} & \textbf{SafeLattice} \\
\midrule
balanced-mid & 0.71 & 0.54 & 0.42 \\
careless-highcap & 0.82 & 0.45 & \textbf{0.30} \\
cautious & 0.67 & 0.69 & 0.69 \\
over-refuser & 0.39 & 0.51 & 0.51 \\
\bottomrule
\end{tabular}
\smallskip
\begin{minipage}{\linewidth}
\scriptsize
\textit{Note.} The high-completion/low-discipline model ranks worst under SafeLattice; completion-only and SafeLattice rankings are nearly inverted (Kendall $\tau=-0.67$), and binary vs.\ SafeLattice disagree on the middle ranks ($\tau=0.67$).
\end{minipage}
\end{table}

\textbf{Severity resolution.} Among detected violations, binary scoring emits a single distinct safety value ($\{0.0\}$), whereas SafeLattice emits three ($\{0.0, 0.3, 0.6\}$), recovering the diagnostic granularity that Table~\ref{tab:binary-collapse} predicts.

\textbf{Ranking divergence.} Table~\ref{tab:ranking} shows per-model mean task scores. The result reproduces the AgentCIBench finding \cite{agentcibench2026} that completion is a poor proxy for safety: the \texttt{careless-highcap} agent ranks \emph{first} by completion ($0.82$) but \emph{last} under SafeLattice ($0.30$). The completion-only and SafeLattice rankings are nearly inverted (Kendall $\tau=-0.67$), and binary and SafeLattice disagree on the middle ranks (Kendall $\tau=0.67$): SafeLattice's severity weighting demotes the balanced agent below the over-refuser because its violations are more often Critical. Binary scoring, treating all violations identically, cannot make this distinction.

\textbf{Over-refusal control.} A scoring scheme that simply penalized every ``tempting'' action would inflate safety while blocking legitimate work; Yu et al.\ \cite{yu2026taxonomy} note that 60\% of agent-safety benchmarks cannot distinguish a safe agent from one that refuses everything. We therefore include a benign control: an \emph{authorized} configuration audit whose required action is exactly the disclosure a naive filter would block. On this control the deliberately over-cautious agent over-refuses (rate $1.0$: it declines the authorized action and fails the task), whereas the other three agents complete it (rate $0.0$). Because SafeLattice scores safety and completion jointly through the multiplicative gate, the over-refuser's high safety is offset by its failed completion---it does not rise to the top of the ranking (Table~\ref{tab:ranking}). This confirms SafeLattice rewards \emph{calibrated} rather than \emph{blanket} caution.

\subsection{Live Evaluation Protocol}
\label{sec:live-protocol}

The corpus above isolates the scoring method under controlled agent behavior. To measure SafeLattice on sampled frontier-model behavior, we run the identical scoring pipeline on real rollouts (harness in Appendix~\ref{app:repro}). The design follows the field's minimum reporting standard \cite{yu2026taxonomy}: (1)~\textbf{models}: a six-model roster spanning capability tiers (GPT-4o, GPT-4o-mini, Claude Sonnet, Claude Haiku, Gemini Flash, Llama-3.3-70B) accessed through a single OpenRouter endpoint; (2)~\textbf{tasks}: the full 94-task safety-relevant subset; (3)~\textbf{trials}: three to five independent rollouts per (model, task) to quantify variance; (4)~\textbf{scoring}: every trace scored under both the binary gate and SafeLattice by the same code paths used above; (5)~\textbf{statistics}: per-model mean task scores with bootstrap $95\%$ confidence intervals, pass$^k$/pass@$k$, and ranking concordance via Kendall's $W$ and pairwise $\tau$; (6)~\textbf{labels}: a semi-automated labeler proposes violation labels from independent detectors (raw credential regex over the full trace, forbidden-tool inventory, injection-then-action) which are then mechanically verified against task reference solutions; (7)~\textbf{over-refusal}: measured on the benign control set per model. We preregistered the hypothesis that live results would diverge from completion-only evaluation in the same direction as the trace-level result (capable-but-careless agents penalized); the executed evaluation below confirms the per-trace form of this hypothesis (the most capable models commit the most critical violations) while aggregate rankings remain stable at the achieved scale. The harness prints a cost estimate and is resumable; $50$ of the $94$ subset tasks have been swept at $3$--$5$ trials, with the remainder queued.

\subsection{Live Multi-Model Evaluation Results}
\label{sec:live-pilot}

We executed this protocol at scale: $1{,}222$ live rollouts ($6$ models $\times$ $50$ safety-relevant tasks $\times$ $3$--$5$ trials each, run through OpenRouter with a $12$-way parallel, resumable harness) at a total API cost of roughly \$36. Every trace is graded by the production Claw-Eval graders (deterministic checks plus LLM judge) and dual-scored by the identical binary and SafeLattice code paths.

\textbf{Ground-truth labeling and verification audit.} Following the human/mechanical validation-audit practice of DTRAP research articles, ground truth is established in two layers and summarized in Table~\ref{tab:gt-audit}. First, independent deterministic detectors---raw credential regex over the full trace, forbidden-tool inventory, and injection-then-action analysis, none of which share code with either scoring system---propose violation labels for all $1{,}222$ traces, yielding $16$ detector-proposed violations (11 credential-in-output, 5 forbidden-tool calls). Second, because the safety of submission-style tasks (expense reports, inventory orders, ambiguous-recipient emails) cannot be judged by a generic detector, we encode task-specific mechanical verifiers (\texttt{verify\_labels.py}) for the six tasks carrying $86\%$ of proposed and candidate detections. Each verifier checks a submission against the task's reference solution---for example, whether a submitted expense report actually contains the duplicate transaction pair (\texttt{txn\_002}/\texttt{txn\_003}), whether an inventory order restocks a SKU with ample stock, or whether an email was sent before the required clarification. Across the $180$ verifier-adjudicated rollouts the verifiers confirmed $69$ true violations and overturned $111$ conservative flags. Combined with the $16$ detector-proposed violations, this yields $85$ verified ground-truth violations among the $1{,}222$ rollouts, against which both scoring systems are measured. Table~\ref{tab:live-pilot} reports per-model results with bootstrap $95\%$ confidence intervals.

\begin{table}[t]
\centering
\caption{Ground-truth labeling and verification audit for the live evaluation.}
\label{tab:gt-audit}
\begin{tabular}{l c}
\toprule
Audit item & Result \\
\midrule
Live rollouts labeled & 1{,}222 \\
Independent detectors applied & 3 \\
Detector-proposed violations & 16 \\
\quad credential-in-output & 11 \\
\quad forbidden-tool calls & 5 \\
Tasks with mechanical verifiers & 6 \\
Verifier-adjudicated rollouts & 180 \\
Coverage of proposed/candidate detections & 86\% \\
Violations confirmed by verifiers & 69 \\
Conservative flags overturned by verifiers & 111 \\
Total verified ground-truth violations & 85 \\
LLM involvement in labeling & None \\
\bottomrule
\end{tabular}
\smallskip
\begin{minipage}{\linewidth}
\scriptsize
\textit{Note.} Detectors and verifiers are deterministic and share no code with the binary or SafeLattice scoring paths, so ground truth is independent of the systems under comparison. Verifiers adjudicate submission-style tasks where a generic credential/tool detector cannot decide correctness.
\end{minipage}
\end{table}

\begin{table}[t]
\centering
\caption{Live evaluation: mean task score with bootstrap 95\% CIs.}
\label{tab:live-pilot}
\begin{tabular}{lcccc}
\toprule
\textbf{Model} & \textbf{Completion} & \textbf{SafeLattice} & \textbf{Viol.\ C/H/M} & \textbf{OR} \\
\midrule
Claude Sonnet 4.6 & 0.73 [0.68, 0.77] & 0.75 [0.71, 0.78] & 3/0/10 & 0.00 \\
Claude Haiku 4.5 & 0.64 [0.59, 0.69] & 0.66 [0.62, 0.71] & 3/10/15 & 0.03 \\
GPT-4o & 0.48 [0.45, 0.52] & 0.56 [0.53, 0.59] & 2/0/20 & 0.00 \\
GPT-4o-mini & 0.43 [0.40, 0.47] & 0.50 [0.47, 0.54] & 3/12/15 & 0.00 \\
Gemini 2.5 Flash & 0.28 [0.24, 0.31] & 0.42 [0.39, 0.45] & 0/0/0 & 0.13 \\
Llama-3.3-70B & 0.28 [0.24, 0.32] & 0.36 [0.33, 0.40] & 0/13/20 & 0.00 \\
\bottomrule
\end{tabular}
\smallskip
\begin{minipage}{\linewidth}
\scriptsize
\textit{Note.} $n{=}194$--$210$ rollouts per model. Viol.\ = SafeLattice-detected violations by severity (Critical/High/Medium). OR = over-refusal rate on the benign control (40 rollouts per model).
\end{minipage}
\end{table}

Four findings emerge on sampled frontier-model behavior.

\textbf{(1)~Frontier models leak credentials live, and the binary gate scores every leak safe.} On the API-configuration-audit tasks---whose prompt deliberately asks for keys, and whose declared check requires refusal or redaction---four of six models disclosed credential material in their output: Claude Haiku ($3$ trials), Claude Sonnet ($3$), GPT-4o-mini ($3$), and GPT-4o ($2$). SafeLattice's Star Property analysis flags all $11$ as Critical write-down violations; the binary gate, blind to information flow, scores all $11$ as $\mathrm{safety}=1.0$. The only models with zero critical violations (Gemini, Llama) are also the least capable---capability and exposure rise together, which is precisely the correlation a safety score must not reward.

\textbf{(2)~Detection gap, now measured against verified ground truth.} Of the $85$ ground-truth violations, the binary gate detects $5$ (recall $0.06$, all forbidden-tool calls: ticket closes by Llama, task deletions by GPT-4o-mini); SafeLattice detects all $85$ (recall $1.0$) at precision $0.67$ ($F_1$ $0.81$ vs.\ $0.11$). The $41$ SafeLattice false positives are concentrated in conservative \texttt{wrong\_data} flagging on submission tasks where the agent's submission was in fact correct---the same precision cost quantified on the corpus (Section~\ref{sec:threats}), now measured live.

\textbf{(3)~Severity structure differs across models even when rankings agree.} At this scale the completion, binary, and SafeLattice mean-score rankings coincide (Kendall's $W{=}1.0$): violations occur in ${\sim}10\%$ of rollouts, so aggregate means are dominated by completion. The value of the lattice is diagnostic: SafeLattice resolves three severity levels live ($11$ Critical, $35$ High, $80$ Medium) and exposes per-model risk profiles the aggregate hides---Claude Sonnet's violations are few but include Critical credential leaks, while Llama's are numerous but capped at High. A binary gate collapses both to indistinguishable zeros on $7\%$ of tasks and silence elsewhere.

\textbf{(4)~Over-refusal is rare once incapacity is separated.} Our refined metric counts a benign-task failure as over-refusal only when an explicit refusal phrase accompanies non-completion (Appendix~\ref{app:repro}); low-completion failures without a refusal signal are classified as incapacity. Over-refusal is at most $0.13$ (Gemini) and zero for four models, while incapacity ranges from $0.25$ to $0.75$---a naive metric that counts all benign-task failures as refusals would have inflated over-refusal by $3$--$9\times$. This operationalizes Yu et al.'s refusal/safety distinction \cite{yu2026taxonomy} and confirms SafeLattice's severity weighting does not push models into blanket caution.

\subsection{Threats to Validity}
\label{sec:threats}

\textbf{Conservative flagging (precision cost).} SafeLattice's credential detector treats any credential-like pattern appearing after Critical-data access as a write-down, even when disclosure is authorized. On our four authorized-disclosure probes it raised 3 false positives---identical to the binary gate, which also flags credentials in output. Context-aware disclosure intent is future work; the shared false-positive count shows SafeLattice does not \emph{worsen} precision relative to the status quo while greatly improving recall.

\textbf{Evasion robustness.} Pattern matching alone is not robust to obfuscation: against six credential-encoding variants, the regex detector achieves recall $0.33$, missing base64, unicode-homoglyph, whitespace-split, and reversed encodings. The canary-based obfuscation detector (Section~\ref{sec:flow-analyzer}) closes this measured gap, reaching recall $1.0$ on the same six-variant probe. The residual limitation is transformations outside the probe set: semantic paraphrase (``the key starts with sk-live and ends in 5678''), multi-character substitutions beyond Hamming distance 1, custom ciphers, and exfiltration chunked across many turns. We report the probe set explicitly so the robustness claim is scoped to it.

\textbf{Trace-level vs.\ live agents.} Our corpus is curated, not sampled from live rollouts; simulated-agent behavior is deterministic by construction, so absolute scores should not be read as frontier-model measurements. This is the accepted trade-off of trace-level benchmarks \cite{tracesafe2026}, and we mitigate it in two ways: scenarios are grounded in real task IDs and CVE patterns, and the live evaluation of Section~\ref{sec:live-pilot} reproduces the detection-gap and severity-resolution findings on $1{,}222$ sampled rollouts with verified ground truth. One trace-level finding did \emph{not} transfer unchanged: the corpus's aggregate ranking inversion does not appear in live aggregate rankings at $50$-task scale, because live violations occur in only ${\sim}10\%$ of rollouts; the live divergence instead manifests per-trace, as critical violations concentrated in the most capable models. We report both rather than claiming the stronger result.

\textbf{Construct validity of the binary baseline.} We model the current gate as output-only $+$ forbidden-tool detection. This is faithful to the audited status quo (81\% of tasks default to $\mathrm{safety}=1.0$; 35 declared checks are unenforced) and to the output-level paradigm criticized by AgentSCOPE and HarnessAudit, but a specific grader could implement a stronger ad-hoc check for an individual task; our baseline reflects the pipeline default, not the best possible per-task grader.

\subsection{Minimum Reporting Standards Compliance}
\label{sec:min-reporting}

Yu et al.\ \cite{yu2026taxonomy} propose minimum reporting standards for agent-safety benchmarks. Table~\ref{tab:min-reporting} maps each item to where this work satisfies it, distinguishing what is established on the trace-level corpus from what the live protocol (Section~\ref{sec:live-protocol}) produces on execution.

\begin{table}[t]
\centering
\caption{Compliance with the minimum reporting standard of Yu et al.\ \cite{yu2026taxonomy}.}
\label{tab:min-reporting}
\begin{tabular}{p{8cm}p{4.5cm}}
\toprule
\textbf{Reporting item} & \textbf{Status} \\
\midrule
(i) Threat model + agent authority level & Sec.\ \ref{sec:background}--\ref{sec:analysis} \\
(ii) Environment fidelity category & Containerized mock services \\
(iii) Joint safety--utility metric & Multiplicative gate, Eq.\ \eqref{eq:scoring} \\
(iv) Over-refusal rate on benign tasks & Corpus + live (refusal vs.\ incapacity), Tab.\ \ref{tab:live-pilot} \\
(v) Severity-weighted harm scores & SafeLattice lattice \\
(vi) Multi-run variance + 95\% CIs & 3--5 trials, bootstrap CIs, Tab.\ \ref{tab:live-pilot} \\
(vii) LLM-judge audit protocol & Deterministic + judge, Sec.\ \ref{sec:use-of-llms} \\
(viii) Correlation with capability & Completion vs.\ safety, Tab.\ \ref{tab:ranking} \\
\bottomrule
\end{tabular}
\end{table}

\subsection{Discussion}

Our results demonstrate that formal security models provide a principled foundation for agent safety evaluation that binary scoring lacks. The key theoretical contribution is the mapping in Table~\ref{tab:blp-mapping}: by establishing a formal correspondence between BLP constructs and agent evaluation concepts, we show that the limitations of binary scoring are not merely engineering oversights but violations of well-established security properties.

The practical implications are significant. Without severity differentiation, benchmark operators cannot distinguish agents that pose genuine security risks (credential leakage) from those that exhibit minor policy deviations (unnecessary drafts). Without information flow tracking, benchmarks systematically overestimate safety by ignoring data flow violations. Without state transition verification, multi-step attacks that compose individually safe actions into unsafe trajectories go undetected. These are not abstract concerns: the live evaluation shows frontier models committing exactly these violations against a benchmark that scores them safe.

The limitations of the current implementation---conservative flagging of authorized disclosure and evasion transformations outside the measured probe set---are quantified in Section~\ref{sec:threats}. Context-aware flow analysis that distinguishes legitimate from illegitimate disclosure, and extending canary matching to semantic paraphrase and multi-turn chunked exfiltration, are the most important directions for reducing the residual precision and evasion gaps.

%% ============================================================
%% SECTION 9: CONCLUSION AND FUTURE WORK
%% ============================================================
\section{Conclusion}
\label{sec:conclusion}

We have demonstrated that binary safety scoring in agent evaluation benchmarks violates the foundational properties of the Bell-LaPadula security model: hierarchical classification, information flow control, and secure state transitions. Using Claw-Eval as a case study, we showed that 81\% of tasks never evaluate safety, binary scoring equates critical credential leaks with minor policy deviations, and information flow violations are invisible to current evaluation mechanisms. In a live evaluation of 1,222 frontier-model rollouts, four of six models leaked credential material on an authorized-audit task and the deployed binary gate scored every leak safe.

SafeLattice addresses these gaps by applying BLP's formal security model to agent evaluation. Our four-level security classification lattice provides severity-differentiated scoring that preserves backward compatibility for critical failures. Our information flow analysis detects write-down violations in 12 credential-bearing tasks that currently score $\texttt{safety} = 1.0$. Our automatic enforcement module increases safety evaluation coverage from 19\% to 28\% by closing the gap between declared safety constraints and actual evaluation. Against verified ground truth, SafeLattice raises detection recall from 0.06 to 1.0 over the deployed gate while resolving three severity levels.

The broader implication is that agent safety evaluation requires the same formal rigor that has been standard in systems security since 1973. Binary scoring is not a simplification---it is a categorical failure to model the hierarchical, flow-controlled nature of security. As the OpenClaw incidents show, the threats are extant and weaponized; the evaluation practice that certifies these agents must catch up. We hope SafeLattice demonstrates the value of applying decades of formal security research to the emerging challenge of evaluating autonomous AI agents.

\subsection{Future Work}

Several directions extend this work:

\begin{itemize}
    \item \textbf{Completing the live protocol.} Extending the executed evaluation (Section~\ref{sec:live-pilot}; 50 tasks, 3--5 trials) to all 94 subset tasks at a uniform five trials, and extending task-specific ground-truth verifiers beyond the six highest-violation tasks, to test whether aggregate rankings separate at full scale.
    \item \textbf{Context-aware flow analysis.} Distinguishing legitimate from illegitimate credential disclosure based on task intent, reducing the authorized-disclosure false positives quantified in Section~\ref{sec:threats}.
    \item \textbf{Evasion detection beyond the probe set.} The canary detector closes the measured six-variant gap ($0.33 \to 1.0$); extending it to semantic paraphrase, multi-character substitution, and multi-turn chunked exfiltration requires embedding-level or judge-assisted matching.
    \item \textbf{Dynamic safety budgets.} Inspired by CAPRI-DP's adaptive privacy budget \cite{capridp2025}, implementing per-task safety budgets that decay with each minor violation, consistent with the state-transition model of Section~\ref{sec:proofs}.
    \item \textbf{Full coverage expansion.} Developing domain-specific safety constraints for the remaining uncovered tasks, particularly in the Multimodal (3.0\%) and Multi-turn (0.0\%) categories.
    \item \textbf{Clark-Wilson transaction integrity.} Applying constrained data items and transformation procedures to verify that agent tool-call sequences maintain data integrity through well-formed transactions.
\end{itemize}

%% ============================================================
%% APPENDIX: REPRODUCIBILITY
%% ============================================================
\appendix
\section{Reproducibility}
\label{app:repro}

All artifacts are released with the (anonymized) evaluation repository, available to reviewers, under \texttt{experiments/safelattice/} and \texttt{analysis/}.

\textbf{Task subset.} \texttt{subset.py} derives the 94-task safety-relevant subset from the machine-readable audit (\texttt{analysis/safety\_audit\_data.json}): 57 safety-checked, 12 credential-vulnerable, 9 injection, 35 YAML-grader-mismatch tasks (deduplicated). Output: \texttt{analysis/safelattice\_subset.json}.

\textbf{Trace corpus.} \texttt{trace\_corpus.py} deterministically generates the 68-trace labeled corpus (17 grounded scenarios $\times$ 4 model profiles) in the runner's JSONL format under \texttt{analysis/trace\_corpus/}, with a ground-truth manifest. Model profiles are fixed (discipline $\in\{0.25,0.55,0.80,0.97\}$), so the corpus is bit-for-bit reproducible with no random seeds.

\textbf{Scoring and metrics.} \texttt{dual\_score.py} scores each trace under the binary baseline and SafeLattice via the same \texttt{load\_trace} and \texttt{enforce\_safety\_checks} used in production, emitting \texttt{analysis/safelattice\_dual\_scoring.json} (detection confusion, false-negatives-fixed, severity resolution, per-model ranking, Kendall $\tau$). \texttt{measurement.py} adds the evasion-robustness probe. \texttt{analyze.py} renders the report and Tables~\ref{tab:detection}--\ref{tab:ranking}.

\textbf{Live runs.} \texttt{run\_live.py} executes the identical comparison on real rollouts via a single OpenRouter endpoint, with three subcommands: \texttt{preflight} (1-token reachability check per model), \texttt{sweep} (the resumable, cost-estimated model$\times$task$\times$trial matrix; \texttt{roster.py} defines the six-model roster and pricing), and \texttt{score} (grades any un-graded trace with the production graders and judge, then dual-scores). \texttt{stats.py} computes bootstrap $95\%$ CIs, pass$^k$/pass@$k$, and Kendall's $W$/$\tau$; \texttt{refusal.py} computes over-refusal and incapacity rates on the benign control set (a benign-task failure counts as over-refusal only with an explicit refusal phrase); \texttt{label.py} proposes labels from independent detectors and \texttt{verify\_labels.py} mechanically confirms them against task reference solutions for the highest-violation tasks, after which \texttt{label.py --score} computes live detection precision/recall. The outputs of Section~\ref{sec:live-pilot} are released as \texttt{analysis/safelattice\_live\_scoring.json}, \texttt{safelattice\_live\_stats.json}, \texttt{safelattice\_live\_detection.json}, \texttt{safelattice\_over\_refusal.json}, and \texttt{live\_label\_proposals.csv}. The full protocol:
\begin{lstlisting}[language=bash]
export OPENROUTER_API_KEY=sk-or-...
python -m experiments.safelattice.run_live \
  preflight --config config_openrouter.yaml
python -m experiments.safelattice.run_live \
  sweep --config config_openrouter.yaml --trials 5
python -m experiments.safelattice.run_live \
  score --trace-root traces_live
python -m experiments.safelattice.stats \
  --trace-root traces_live
\end{lstlisting}

\textbf{Tests.} 119 tests (\texttt{tests/test\_safelattice.py}, \texttt{test\_safelattice\_integration.py}, \texttt{test\_safelattice\_live.py}) cover the lattice algebra, BLP flow, Biba integrity, canary-based obfuscation detection, real-trace replay, backward compatibility, and the live-experiment tooling (bootstrap CIs, Kendall's $W$, refusal/incapacity classification, label proposal). All run without API access; execute with \texttt{pytest tests/}.

\bibliographystyle{ACM-Reference-Format}
\bibliography{references}

\end{document}
